% generated by tools/compass.py
\begin{tikzpicture}[x=1mm,y=1mm,font=\sffamily]
\fill[accentlight,rounded corners=2pt] (0,-69.1) rectangle (136,7.1);
\node[anchor=north west,font=\sffamily\bfseries\footnotesize,text=accent] at (1.5,5.6) {Vol.\,1};
\node[anchor=north west,font=\sffamily\scriptsize,text=quiet,align=left] at (1.5,1.3) {Numbers and\\Equations};
\fill[accentlight,rounded corners=2pt] (0,-109.3) rectangle (136,-70.3);
\node[anchor=north west,font=\sffamily\bfseries\footnotesize,text=accent] at (1.5,-71.8) {Vol.\,2};
\node[anchor=north west,font=\sffamily\scriptsize,text=quiet,align=left] at (1.5,-76.1) {Algebra and\\Rank};
\fill[accentlight,rounded corners=2pt] (0,-149.5) rectangle (136,-110.5);
\node[anchor=north west,font=\sffamily\bfseries\footnotesize,text=accent] at (1.5,-112.0) {Vol.\,3};
\node[anchor=north west,font=\sffamily\scriptsize,text=quiet,align=left] at (1.5,-116.3) {States and\\Spacetime};
\fill[accentlight,rounded corners=2pt] (0,-177.3) rectangle (136,-150.7);
\node[anchor=north west,font=\sffamily\bfseries\footnotesize,text=accent] at (1.5,-152.2) {Vol.\,4};
\node[anchor=north west,font=\sffamily\scriptsize,text=quiet,align=left] at (1.5,-156.5) {Action and\\Dynamics};
\draw[draw=black!55,line width=0.5pt,fill=white,rounded corners=2pt] (25,-5.3) rectangle (135,5.3);
\node[anchor=west,font=\sffamily\bfseries\scriptsize] at (26.8,2.3) {1.\ Why could Hamilton build numbers in four dimensions but not in three?};
\node[anchor=west,font=\sffamily\scriptsize,text=quiet] at (26.8,-2.4) {$\to$ multiplication that keeps length: $q\bar q=|q|^2$; only dimensions 1, 2, 4, 8};
\node[anchor=east,font=\sffamily\tiny,text=accent] at (133.8,2.3) {Part I};
\draw[-{Stealth[length=1.6mm]},black!55] (80.0,-5.3) -- (80.0,-7.1);
\draw[draw=black!55,line width=0.5pt,fill=white,rounded corners=2pt] (25,-17.7) rectangle (135,-7.1);
\node[anchor=west,font=\sffamily\bfseries\scriptsize] at (26.8,-10.1) {2.\ What does a unit quaternion do?};
\node[anchor=west,font=\sffamily\scriptsize,text=quiet] at (26.8,-14.8) {$\to$ rotations of space (covered twice), the sphere of imaginary units, Hopf circles};
\node[anchor=east,font=\sffamily\tiny,text=accent] at (133.8,-10.1) {Part I};
\draw[-{Stealth[length=1.6mm]},black!55] (80.0,-17.7) -- (80.0,-19.5);
\draw[draw=black!55,line width=0.5pt,fill=white,rounded corners=2pt] (25,-30.1) rectangle (135,-19.5);
\node[anchor=west,font=\sffamily\bfseries\scriptsize] at (26.8,-22.5) {3.\ How do we compute with numbers that do not commute?};
\node[anchor=west,font=\sffamily\scriptsize,text=quiet] at (26.8,-27.2) {$\to$ left and right matrices; complexifying brings nonzero elements of length zero};
\node[anchor=east,font=\sffamily\tiny,text=accent] at (133.8,-22.5) {Part I};
\draw[-{Stealth[length=1.6mm]},black!55] (80.0,-30.1) -- (80.0,-31.9);
\draw[draw=black!55,line width=0.5pt,fill=white,rounded corners=2pt] (25,-42.5) rectangle (135,-31.9);
\node[anchor=west,font=\sffamily\bfseries\scriptsize] at (26.8,-34.9) {4.\ Can complex analysis be carried to four dimensions?};
\node[anchor=west,font=\sffamily\scriptsize,text=quiet] at (26.8,-39.6) {$\to$ the Fueter operator $D$, with $\bar D D=\Delta$};
\node[anchor=east,font=\sffamily\tiny,text=accent] at (133.8,-34.9) {Part II};
\draw[-{Stealth[length=1.6mm]},black!55] (80.0,-42.5) -- (80.0,-44.3);
\draw[draw=black!55,line width=0.5pt,fill=white,rounded corners=2pt] (25,-54.9) rectangle (135,-44.3);
\node[anchor=west,font=\sffamily\bfseries\scriptsize] at (26.8,-47.3) {5.\ What happens if the coefficients $1,i,j,k$ are cycled?};
\node[anchor=west,font=\sffamily\scriptsize,text=quiet] at (26.8,-52.0) {$\to$ a common linear solution appears: the coefficient matrix $C$ is singular (rank 3)};
\node[anchor=east,font=\sffamily\tiny,text=accent] at (133.8,-47.3) {Part II};
\draw[-{Stealth[length=1.6mm]},black!55] (80.0,-54.9) -- (80.0,-56.7);
\draw[draw=black!55,line width=0.5pt,fill=white,rounded corners=2pt] (25,-67.3) rectangle (135,-56.7);
\node[anchor=west,font=\sffamily\bfseries\scriptsize] at (26.8,-59.7) {6.\ Is that singularity an accident?};
\node[anchor=west,font=\sffamily\scriptsize,text=quiet] at (26.8,-64.4) {$\to$ no: the DFT eigenvalues are paravectors of length zero (nullity theorem)};
\node[anchor=east,font=\sffamily\tiny,text=accent] at (133.8,-59.7) {Part III};
\draw[-{Stealth[length=1.6mm]},black!55] (80.0,-67.3) -- (80.0,-72.1);
\draw[draw=black!55,line width=0.5pt,fill=white,rounded corners=2pt] (25,-82.7) rectangle (135,-72.1);
\node[anchor=west,font=\sffamily\bfseries\scriptsize] at (26.8,-75.1) {7.\ Which algebra hides in the orderings and coefficients?};
\node[anchor=west,font=\sffamily\scriptsize,text=quiet] at (26.8,-79.8) {$\to$ $S_4$ and the Klein group, Clifford systems, octonions};
\node[anchor=east,font=\sffamily\tiny,text=accent] at (133.8,-75.1) {Part IV};
\draw[-{Stealth[length=1.6mm]},black!55] (80.0,-82.7) -- (80.0,-84.5);
\draw[draw=accent,line width=1.1pt,fill=accent!12,rounded corners=2pt] (25,-95.1) rectangle (135,-84.5);
\node[anchor=west,font=\sffamily\bfseries\scriptsize] at (26.8,-87.5) {8.\ What shape does the set of singular configurations have?};
\node[anchor=west,font=\sffamily\scriptsize,text=quiet] at (26.8,-92.2) {$\to$ rank geometry: two mechanisms by handedness; odd rank 7 after complexifying};
\node[anchor=east,font=\sffamily\tiny,text=accent] at (133.8,-87.5) {Part V};
\draw[-{Stealth[length=1.6mm]},black!55] (80.0,-95.1) -- (80.0,-96.9);
\draw[draw=black!55,line width=0.5pt,fill=white,rounded corners=2pt] (25,-107.5) rectangle (135,-96.9);
\node[anchor=west,font=\sffamily\bfseries\scriptsize] at (26.8,-99.9) {9.\ Why does one formula have two real forms?};
\node[anchor=west,font=\sffamily\scriptsize,text=quiet] at (26.8,-104.6) {$\to$ $t^2+|\mathbf x|^2$ and $t^2-|\mathbf x|^2$; on the Lorentz side rank 7 is a shared light ray};
\node[anchor=east,font=\sffamily\tiny,text=accent] at (133.8,-99.9) {Part VI};
\draw[-{Stealth[length=1.6mm]},black!55] (80.0,-107.5) -- (80.0,-112.3);
\draw[draw=black!55,line width=0.5pt,fill=white,rounded corners=2pt] (25,-122.9) rectangle (135,-112.3);
\node[anchor=west,font=\sffamily\bfseries\scriptsize] at (26.8,-115.3) {10.\ What does length zero mean in quantum theory?};
\node[anchor=west,font=\sffamily\scriptsize,text=quiet] at (26.8,-120.0) {$\to$ $\det\rho=0$: pure states; entanglement of two qubits};
\node[anchor=east,font=\sffamily\tiny,text=accent] at (133.8,-115.3) {Part VII};
\draw[-{Stealth[length=1.6mm]},black!55] (80.0,-122.9) -- (80.0,-124.7);
\draw[draw=black!55,line width=0.5pt,fill=white,rounded corners=2pt] (25,-135.3) rectangle (135,-124.7);
\node[anchor=west,font=\sffamily\bfseries\scriptsize] at (26.8,-127.7) {11.\ What does length zero mean in relativity?};
\node[anchor=west,font=\sffamily\scriptsize,text=quiet] at (26.8,-132.4) {$\to$ light rays = rank one = spinors; mass is measured from the equator};
\node[anchor=east,font=\sffamily\tiny,text=accent] at (133.8,-127.7) {Part VIII};
\draw[-{Stealth[length=1.6mm]},black!55] (80.0,-135.3) -- (80.0,-137.1);
\draw[draw=black!55,line width=0.5pt,fill=white,rounded corners=2pt] (25,-147.7) rectangle (135,-137.1);
\node[anchor=west,font=\sffamily\bfseries\scriptsize] at (26.8,-140.1) {12.\ What survives when states and spacetime move?};
\node[anchor=west,font=\sffamily\scriptsize,text=quiet] at (26.8,-144.8) {$\to$ Berry phase, Thomas--Wigner rotation, curved space};
\node[anchor=east,font=\sffamily\tiny,text=accent] at (133.8,-140.1) {Part IX};
\draw[-{Stealth[length=1.6mm]},black!55] (80.0,-147.7) -- (80.0,-152.5);
\draw[draw=black!55,line width=0.5pt,fill=white,rounded corners=2pt] (25,-163.1) rectangle (135,-152.5);
\node[anchor=west,font=\sffamily\bfseries\scriptsize] at (26.8,-155.5) {13.\ Kinematics is complete. What drives it?};
\node[anchor=west,font=\sffamily\scriptsize,text=quiet] at (26.8,-160.2) {$\to$ the action principle: Maxwell, gravity, throats, collapse, quantisation};
\node[anchor=east,font=\sffamily\tiny,text=accent] at (133.8,-155.5) {Part X--XI};
\draw[-{Stealth[length=1.6mm]},black!55] (80.0,-163.1) -- (80.0,-164.9);
\draw[draw=black!55,line width=0.5pt,fill=white,rounded corners=2pt] (25,-175.5) rectangle (135,-164.9);
\node[anchor=west,font=\sffamily\bfseries\scriptsize] at (26.8,-167.9) {14.\ What do we still not know?};
\node[anchor=west,font=\sffamily\scriptsize,text=quiet] at (26.8,-172.6) {$\to$ open problems; the fine-structure constant $\alpha$ is not derived};
\node[anchor=east,font=\sffamily\tiny,text=accent] at (133.8,-167.9) {Part Epilogue};
\end{tikzpicture}
